\section{Idea 3: assign flexibility to distinct obstructions}\label{sec:gap}
\intuition{Deleting a freely movable alternative can merge many orders.
The Condorcet restrictions force a compensating rigid adjacent pair
elsewhere. The key is to make this assignment injective.}

\begin{definition}[Moves and their counts]\label{def:moves}
A lift of \(x\) moves it upwards, preserving the relative order of all
other labels and staying inside \(D\); a drop moves it downwards.
Moves may jump several positions. An incidence \((\sigma,x)\), with
\(\sigma\in D\), is free if both moves exist and frozen if neither
exists. Their counts are \(\fr,\fz\). A neighbouring ordered pair
\(a,c\) in \(\sigma\) is locked if \(c\) has no lift and \(a\)
has no drop. Let \(\lk\) count these order/pair incidences.
\end{definition}
\lean{\path{Condorcet.HasLift}, \path{HasDrop}, \path{Free},
\path{Frozen}, \path{freeCount}, \path{frozenCount}, \path{Locked},
\path{lockedCount}. The exact exported move names and their source
locations are recorded in the definition index. Here locked refers
to arbitrary lifts and drops, not only the adjacent swap.}

\begin{lemma}[Two move counts]\label{lem:moves}
Every family satisfies
\[
 2\Sig=nM-\fr+\fz.
\]
For a full domain one also has
\[
 2\Sig=(n+1)M-A+\lk,
 \qquad
 3\Sig+\fr=(n+1)M+G+\fz+\lk.             \tag{\(*\)}
\]
\end{lemma}
\begin{proof}
Each nonempty fibre has exactly one highest occupied slot and one
lowest. Thus \(\Sig\) counts incidences with no lift, and also
incidences with no drop. An incidence contributes \(0,1,2\) to their
sum according as it is free, has just one move, or is frozen. This
gives the first equation.

For a neighbouring pair \(a,c\), count the missing drop of \(a\)
and missing lift of \(c\). An enabled swap contributes zero, a locked
pair contributes two. Every other blocked swap contributes one.
To justify the last assertion, a blocked swap violates a defining
triple law involving both \(a,c\), since all other comparisons are
unchanged. If the third label \(z\) lies above both, the changed triple
is \(zac\to zca\); its only possible newly violated laws are
\(aN_3\) and \(cN_2\). The first rules out every lift of \(c\)
across \(a\); the second rules out every drop of \(a\) across \(c\).
If \(z\) lies below both, the changed triple is \(acz\to caz\),
and the possible laws are \(aN_2,cN_1\); these likewise rule out one
of the crossing moves. Therefore a forbidden swap cannot have both
crossing moves available. Fullness is what supplies a violated
chosen law when the swapped order is absent.

There are also two end contributions per order: its first label has
no lift and its last label no drop. Summing gives
\(2\Sig=2M+(n-1)M-A+\lk\). Combine this with the first equation
and \(\Sig+G+A=nM\) to obtain \((*)\).
\end{proof}
\lean{\path{Condorcet.first_move_count},
\path{IsFull.missing_swap}, \path{IsFull.second_move_count},
\path{IsFull.gap_slack_identity}. The last is written without
subtraction in \(\mathbb N\):
\(3\,\texttt{deletionSum}+\texttt{freeCount}
=(k+2)|D|+\texttt{gapCount}+\texttt{frozenCount}+\texttt{lockedCount}\).
The coefficient \(k+2\) is \(n+1\) because \(n=k+1\).}

\begin{lemma}[A fixed pair cannot be consecutive in a maximal domain]\label{lem:fixedpair}
If every order in a maximal Condorcet domain has \(a>c\), then no
order in the domain puts \(a,c\) consecutively.
\end{lemma}
\begin{proof}
Otherwise interchange that consecutive pair to obtain \(\tau\).
Triples not containing both labels are unchanged from the original
order. On a triple containing \(a,c,z\), all old orders satisfy
\(cN_1\) and \(aN_3\), because \(a>c\) is fixed. If \(z\) lies
above the consecutive pair in the original order, \(\tau\) has trace
\(zca\), which still satisfies \(cN_1\). If \(z\) lies below,
\(\tau\) has trace \(caz\), still satisfying \(aN_3\). Thus the
old family together with \(\tau\) has a common law on every triple
and is Condorcet. The new order reverses the fixed pair and was absent,
contradicting maximality.
\end{proof}
\lean{\path{Condorcet.IsMaximal.fixed_pair_not_consecutive}.}

\begin{proposition}[Free-to-locked injection]\label{prop:injection}
For a maximal Condorcet domain, \(\fr\leq\lk\).
\end{proposition}
\begin{proof}
Let \(b\) be free in an order whose immediate neighbours are \(a,b,c\).
The order itself, a lift and a drop show the triple traces
\(abc,bac,acb\). The only possible common laws are \(aN_3\) and
\(cN_1\). Both cannot hold: together they force \(a>c\) in every
order, while a lift of \(b\) puts \(a,c\) consecutively, contrary
to Lemma~\ref{lem:fixedpair}.

If \(aN_3\) holds, take the lowest available lift of \(b\); otherwise
take its highest available drop. In the resulting order \(\tau\),
\(a,c\) are consecutive and locked. In the lift case, \(b\) lies
above both: lifting \(c\) above \(a\), or dropping \(a\) below
\(c\), makes \(a\) last on their triple, violating \(aN_3\).
The drop case is reversed. This defines the assignment.

For a fixed moved label \(b\), its preimage is uniquely recovered
by moving \(b\) into the slot between the image pair. It remains to
exclude different moved labels with the same image. The following
three cases give that exclusion explicitly.

Call the image pair \(p,q\). A preimage is upper when its moved label
\(b\) is above \(p\) in \(\tau\). Then \(pN_3\) holds and
\(qN_1\) does not, so there exists \(\chi_b\in D\) with trace
\(qpb\). For a lower preimage \(qN_1\) holds and \(pN_3\) does not.
Suppose two distinct labels \(b_1,b_2\) give the same image.

\emph{Both upper.} Drop \(b_i\) in its original preimage \(\sigma_i\),
obtaining \(\rho_i\). On \(\{b_1,b_2,q\}\) the four traces are
\[
 \begin{array}{c|cccc}
  \text{order}&\sigma_1&\sigma_2&\rho_1&\rho_2\\\hline
  \text{trace}&b_2b_1q&b_1b_2q&b_2qb_1&b_1qb_2.
 \end{array}
\]
Both \(b_1,b_2\) occupy every position, so only \(qN_1\) can hold.
In \(\chi_{b_1}\), \(q\) is above \(p,b_1\). Thus \(b_2\)
must be above \(q\), giving trace \(b_2qp\) on \(\{p,q,b_2\}\),
contrary to its \(pN_3\) law.

\emph{Both lower.} Reverse every order. Maximality and the Condorcet
property are preserved; lifts and drops exchange roles. The assignment
also reverses: exactly one of its two possible triple laws holds, and
the lowest lift becomes the highest drop. This reduces to the case
just excluded.

\emph{One upper and one lower.} Let \(b_1\) be upper and \(b_2\)
lower. The traces on \(\{p,q,b_1,b_2\}\) are
\[
 \tau:b_1pqb_2,\quad\sigma_1:pb_1qb_2,\quad\sigma_2:b_1pb_2q.
\]
Take \(\chi=\chi_{b_1}\), with trace \(qpb_1\). The law \(qN_1\)
on \(\{p,q,b_2\}\) forces \(b_2\) above \(q\), so \(\chi\)
has trace \(b_2qpb_1\). On \(\{b_1,b_2,q\}\), the orders
\(\tau,\sigma_2,\chi\) have traces
\(b_1qb_2,b_1b_2q,b_2qb_1\). These permit just two common laws:
\(b_1N_2\) and \(qN_1\). Let \(\rho\) be a drop of \(b_1\)
from \(\sigma_1\). In \(\rho\), the label \(q\) is above both
\(b_1\) and \(b_2\): the drop puts \(b_1\) below its neighbour
\(q\), while the comparison of \(q\) with \(b_2\) stays fixed.
Thus \(\rho\) rules out \(qN_1\), leaving \(b_1N_2\).
Let \(\lambda\) be a lift of \(b_2\) from \(\sigma_2\).
The remaining law forces the traces of \(\rho,\lambda\) on this
triple to be \(qb_2b_1\) and \(b_1b_2q\).
On \(\{p,b_1,b_2\}\), the five orders
\(\tau,\sigma_1,\rho,\lambda,\chi\) therefore give
\[
 b_1pb_2,\quad pb_1b_2,\quad pb_2b_1,\quad b_1b_2p,\quad b_2pb_1.
\]
Every alternative occupies every position, leaving no common law.
This contradiction excludes the last collision and proves injectivity.
\end{proof}
\lean{\path{CDCert.FreeLocked.FreePair.law}, \path{not_both},
\path{lift_locked}, \path{drop_locked}, followed by
\path{CDCert.FreeLocked.psi_inj} and \path{free_le_locked}.
The conclusion assumes only \path{IsMaximal}; fullness and the triple
laws are obtained inside the dependency chain.}

\begin{theorem}[Gap inequality]\label{thm:gap}
For every maximal Condorcet domain,
\[
 \boxed{G+(n+1)M\leq3\Sig.}
\]
\end{theorem}
\begin{proof}
Use \(\fr\leq\lk\) in Lemma~\ref{lem:moves}, and discard the
nonnegative frozen count. No finite enumeration is used.
\end{proof}
\lean{\path{CDCert.Recurrences.gap_inequality} and its final wrapper
\path{CDCert.TwoInputs.first_input}. The helper
\path{CDCert.Bounds.FreeLockedInjection} is an interface, not an
unproved premise: \path{CDCert.Recurrences.freeLockedInjection}
discharges it using the injection just proved.}
